%======================================================================
%  06 总结提升
%======================================================================
\section{总结提升}

\subsection{一表通览}

\begin{center}
\small
\begin{tabularx}{\linewidth}{>{\sffamily\bfseries\centering\arraybackslash}m{1.6cm} Y Y Y}
\toprule
\textbf{现象} & \multicolumn{1}{c}{\textbf{本质与条件}} & \multicolumn{1}{c}{\textbf{关键特点}} & \multicolumn{1}{c}{\textbf{典型应用}} \\
\midrule
感生电场 & 变化的磁场在周围空间激发电场，与有无导体无关 & 电场线闭合；做功与路径有关；是感生电动势的非静电力 & 电子感应加速器 \\
涡流 & 整块导体内部回路的磁通量变化 & 电阻小、电流大；有热效应、磁效应、力效应；$f$ 越高越强 & 冶炼炉、电磁炉、金属探测器；硅钢片叠压防涡流 \\
电磁阻尼 & 导体相对磁场运动（磁场静止） & 阻碍导体运动；速度越大阻力越大，不能使物体静止；机械能 $\to$ 内能 & 仪表阻尼、电磁刹车、精密仪器减振 \\
电磁驱动 & 磁场相对导体运动 & 导体跟随磁场同向运动，速度小于磁场速度（异步） & 感应电动机、车速表、感应式电能表 \\
\bottomrule
\end{tabularx}
\end{center}

\subsection{楞次定律中“阻碍”的四种表现}

\begin{center}
\small
\begin{tabularx}{\linewidth}{>{\centering\arraybackslash}m{1.8cm} >{\centering\arraybackslash}m{3.6cm} Y}
\toprule
\textbf{口诀} & \textbf{阻碍的对象} & \multicolumn{1}{c}{\textbf{在本讲中的体现}} \\
\midrule
增反减同 & 阻碍磁通量的变化 & 判断涡流方向：线圈电流增大时，涡流与线圈电流方向相反 \\
来拒去留 & 阻碍导体与磁场的相对运动 & 电磁阻尼（阻碍导体运动）、电磁驱动（导体跟随磁场运动）；磁体在铜管中下落变慢 \\
增缩减扩 & 通过改变回路面积阻碍磁通量变化 & 磁感线单向穿过闭合金属环时，磁场增强，环有收缩的趋势 \\
增离减靠 & 通过改变距离阻碍磁通量变化 & 跳环实验：闭合开关瞬间，套在铁芯上的闭合铝环跳起 \\
\bottomrule
\end{tabularx}
\end{center}

\subsection{解题思维流程}

\begin{fangfa}[涡流类问题“四步走”]
\begin{steps}
  \item \textbf{找“变”}：穿过导体（内部回路）的磁通量为什么变化？是磁场在变（感生），还是导体与磁场有相对运动（动生）？
  \item \textbf{判“流”}：有没有闭合回路（整块导体、闭合线圈）？用楞次定律（增反减同）判断涡流或感应电流的方向。
  \item \textbf{析“效”}：分析涡流的热效应（焦耳热、能量转化）、磁效应（对原线圈的影响）或力效应（安培力阻碍相对运动：阻尼还是驱动）。
  \item \textbf{算“量”}：定量问题建立导体棒或线框模型，用动力学观点分析过程，用动量定理求电荷量、位移、时间，用能量守恒求焦耳热。
\end{steps}
\end{fangfa}

\subsection{易错辨析}

\noindent 判断下列说法是否正确，正确的打“\dui”，错误的打“\cuo”。
\begin{enumerate}[label=（\arabic*）,leftmargin=2.6em,itemsep=2pt]
  \item 感生电场是由变化的磁场激发的，与空间中是否存在闭合回路无关。\pd{\dui}{变化的磁场在周围空间激发感生电场，回路只是显示它的工具。}
  \item 感生电场的电场线是闭合的，所以感生电场力对电荷做功与路径无关。\pd{\cuo}{感生电场力做功与路径有关，沿闭合路径一周做功一般不为零。}
  \item 圆柱形区域内的磁场均匀变化时，区域外部没有磁场，所以也没有感生电场。\pd{\cuo}{区域外部同样存在感生电场，$E_{\text{感}}=\dfrac{kR^2}{2r}$。}
  \item 静止的金属块放在磁感应强度很大的恒定磁场中，也会产生很强的涡流。\pd{\cuo}{恒定磁场中静止的金属块，磁通量不变，不产生涡流。}
  \item 涡流是在导体内部自成回路的感应电流，同样遵循楞次定律。\pd{\dui}{涡流本质上是感应电流。}
  \item 在其他条件不变时，交变电流的频率越高，金属中的涡流越大。\pd{\dui}{频率越高，磁通量变化率越大，感应电动势越大，涡流越大。}
  \item 电磁炉面板发热，是因为面板中产生了涡流。\pd{\cuo}{面板是绝缘体，不产生涡流，热量来自锅底。}
  \item 电磁炉和微波炉都是利用涡流加热食物的。\pd{\cuo}{微波炉利用微波使食物中的水分子振动加热，与涡流无关。}
  \item 变压器的铁芯用硅钢片叠成，主要是为了增强磁场。\pd{\cuo}{增强磁场是铁芯的作用，用绝缘硅钢片叠成是为了减小涡流损耗。}
  \item 金属探测器可以检测出藏在行李中的塑料制品。\pd{\cuo}{塑料是绝缘体，不能产生涡流。}
  \item 电磁阻尼中，安培力总是阻碍导体运动，因此能使运动的导体最终保持静止在任意位置。\pd{\cuo}{速度为零时安培力为零，电磁阻尼只能减速，不能“锁住”物体。}
  \item 强磁体在竖直铜管中由静止下落，若铜管足够长，磁体最终做匀速运动，此时减少的重力势能全部转化为内能。\pd{\dui}{匀速时动能不变，重力做功全部转化为内能。}
  \item 电磁驱动中，稳定运动时导体的转速可以等于磁场的转速。\pd{\cuo}{转速相等时无相对运动，无感应电流和驱动力，导体转速必小于磁场转速。}
  \item 阿拉果圆盘实验中，磁针转动是因为圆盘中的自由电子随圆盘转动形成了电流。\pd{\cuo}{圆盘整体呈电中性，随圆盘转动不形成电流；使磁针转动的是涡流的磁场。}
\end{enumerate}
